
en
posant pour tout \(L\in\Pic(X)\)
\[
 \delta(L)=1-L^{-1}.
\]

\textbf{Remarques 5.2.2.} :
\begin{enumerate}[label=\alph*)]
\item Supposons que \(X\) soit le topos localement annelé des faisceaux zariskiens sur un schéma \(S\). Lorsque \(D\) est un diviseur de Cartier \(\geq0\) sur \(X\) et
\[
 L=\cl^{\bullet}(\underline{O}_S(D)),
\]
alors on a
\begin{equation}\tag{5.2.3}
 \theta_S(\delta(L))=\cl_\bullet(\underline{O}_D)
\end{equation}
où \(\theta_S:K^{\bullet}(X)\to K_\bullet(S)\) est l'homomorphisme canonique de 1.2.3. Cela résulte immédiatement de la suite exacte
\[
0\longrightarrow\underline{O}_S(-D)\longrightarrow\underline{O}_S\longrightarrow\underline{O}_D\longrightarrow0.
\]

\item Sous les hypothèses de 5.1 et 5.2 on a pour tout \(L\in\Pic(X)\)
\begin{equation}\tag{5.2.4}
 \delta(L)\equiv L-1\quad\operatorname{mod.}\Filt^2(X),
\end{equation}
ce qui justifie la notation \(c_1(L)\) introduite plus bas. En effet, on a
\[
 (L-1)-(1-L^{-1})=(L-1)(1-L^{-1})\in\Filt^2(X).
\]
\end{enumerate}

\subsection*{5.3.}
La composition de \(\delta\) avec l'épimorphisme canonique
\[
 \Filt^1(X)\longrightarrow\Gr^1(X)
\]
donne la «~première classe de Chern~» sur \(\Pic(X)\) :
\begin{equation}\tag{5.3.1}
 c_1:\Pic(X)\longrightarrow\Gr^1(X).
\end{equation}

